% Side Channels in Filtered Retrieval: body of Section~\ref{sec:side} (the \section heading and its
% label stay in main.tex). Labels defined here: sec:side-game, sec:side-setup, tab:side-hyp,
% sec:side-results. Source: benchmarks/paper/PREREG_X3_X6.md on branch paper/x3-x6 (X6, with
% amendments 1 to 3, all made before any full result). Results marked \tbd{} come from the full run.

Proposition~\ref{prop:reduction} bounds what a subject can \emph{read}; it says nothing about what
the timing or shape of an answer reveals. B\"{u}ttcher and Clarke required of secure full-text
search that no aspect of the output, its order included, depend on content the user may not
search~\cite{buttcher2005security}. Each enforcement path of Section~\ref{sec:system} departs from
that requirement in a predictable way: a widening post-filter works longer when unreadable chunks
crowd the top of the ranking, an approximate filtered search passes unreadable nodes on its way to
readable ones, and a keyword score computed over the whole collection counts unreadable documents.
We ask how much an authenticated subject can learn from these departures about whether unreadable
chunks exist near a query, and what removing each one costs.

\subsection{A distinguishing game}\label{sec:side-game}

Two worlds of equal size differ only in where $m$ unreadable chunks lie. In $W_1$ they are the $m$
passages most similar to a probe query; in $W_0$ each is replaced by a passage of another domain, far
from every probe and cut to the same length, at the same position in the index. The attacker's
readable set is identical in both worlds. It issues the probe and sees only what a caller of the
layer sees: the latency, the number of hits, the completeness flag $c$, the mean score of the hits,
and the hits themselves, from which it can compute recall@10 against the exact top ten of its own
readable set, since it can read and embed that set. The \emph{differencing attacker}, our primary
adversary, sees the same probe in both worlds, as someone who repeats a query before and after a
document is added; its success is the paired AUC, the share of probes whose $W_1$ observation lies
further in the predicted direction than their $W_0$ observation, and its advantage is $2\,\mathrm{AUC}
- 1$. The \emph{one-shot attacker} sees one observation and must threshold it; its success is the AUC
over all $W_1$ against all $W_0$ observations.

\subsection{Setup}\label{sec:side-setup}

\paragraph{Data and probes.}
The background is all 57,638 FiQA passages~\cite{maia2018www18}, with its 6,648 queries as probe
candidates; far chunks come from 12,000 SciDocs abstracts, both from BEIR~\cite{thakur2021beir}, and
all text is embedded with all-MiniLM-L6-v2 (384 dimensions)~\cite{reimers2019sentencebert,wang2020minilm}.
Each of five seeds places eight groups of nine probes, one with $m = 1{,}000$ secrets, two with 100,
three with 10 and three with 1. A probe's secrets are its $m$ nearest passages, excluding hubs, the 1\% of
passages found among the 2,000 nearest neighbors of the largest number of queries; within a group, no probe's secrets fall in the
zone a group mate can explore, its nearest passages up to 1,000 beyond its own secrets, which is the
post-filter's cap. One index $W_1$ holds every group's secrets; for each group $g$, an index $W_0^g$
replaces only group $g$'s secrets by far chunks, so a probe of group $g$ compares two worlds that
differ in its own secrets alone. Background passages fall into 100 permission buckets, and three
attackers read 1\%, 10\% and 50\% of them; secrets carry an ACL none of them holds. Every seed
builds every index afresh.

\paragraph{Configurations.}
The \emph{pre-filter} runs each store's approximate filtered search at documented defaults:
pgvector, Qdrant and Elasticsearch, with FAISS over a flat index and \sys{}'s in-memory backend as
exact controls. The \emph{post-filter} is \sys{}'s widening loop at its defaults: fetch $4k$, widen
fourfold until $k$ hits are readable, the store runs out, or 1,000 candidates have been examined.
Four mitigations are measured against them. \emph{Constant-work widening} (M1) fetches the cap of
1,000 in one round, so every request examines the same number of candidates. \emph{Exact
pre-filtered search} (M2) ranks every readable chunk: pgvector with index scans disabled, Qdrant's
exact flag, Elasticsearch's \texttt{script\_score}. \emph{Separation} (M3) gives the attacker an index
that holds only its own tenant's chunks, as a tenant-per-index layout would; its advantage is zero by
construction, and measuring it twice provides an A/A null for the harness. \emph{Latency padding}
(M4) delays each answer to the next multiple of a budget fixed on the first seed and is evaluated on
the other four; it models an ideal padder and ignores scheduler jitter. For collection statistics,
keyword search runs under the pre-filter on Elasticsearch (BM25), on Qdrant (sparse vectors with an
IDF modifier) and on pgvector (\texttt{ts\_rank\_cd}, which uses no collection statistics: a
negative control), and Elasticsearch adds reciprocal-rank fusion of vector and keyword search.

\paragraph{Measurement and statistics.}
For every seed, store, probe, attacker and configuration, one warm-up is followed by five timed
repetitions, each of which runs $W_0^g$, $W_1$ and the two separated indexes in a fresh random order,
so drift affects every world alike; an observation is the median over repetitions. Every retrieval is
re-checked, and a hit outside the attacker's readable set would stop the run. Intervals are 95\%
cluster-bootstrap intervals over (seed, probe) pairs with 10,000 resamples. A channel is
\emph{present} when the paired AUC reaches the hypothesis's threshold and its interval's lower end
exceeds 0.55, and \emph{absent} when the interval lies within $[0.40, 0.60]$. Two guards precede any
verdict: the A/A comparison of the two separated indexes must show no latency channel, and exact
configurations must return identical hit lists in $W_0^g$ and $W_1$ for at least 99\% of probes. The
experiment runs on GitHub-hosted runners with four virtual CPUs and 16\,GB of memory.

\begin{table}[tbp]
\centering
\small
\begin{tabular}{@{}p{0.27\linewidth}p{0.47\linewidth}p{0.16\linewidth}@{}}
\toprule
Hypothesis & Prediction & Verdict \\
\midrule
H1 widening & Post-filter: count AUC $\ge 0.95$ at $m = 1{,}000$; latency AUC $\ge 0.75$ at $m = 100$
  ($p = 10\%, 50\%$); no channel at $m = 1$ & \tbd{} \\
H2 constant work (M1) & No latency channel; the count channel remains at $m = 1{,}000$ & \tbd{} \\
H3 approximate pre-filter & pgvector and Elasticsearch: latency AUC $\ge 0.75$ at $m = 1{,}000$;
  Qdrant: no channel at $p = 1\%, 10\%$ & \tbd{} \\
H4 exact pre-filter (M2) & No channel in any observable & \tbd{} \\
H5 collection statistics & BM25 and IDF-weighted sparse: score AUC $\ge 0.90$ at every $m$;
  \texttt{ts\_rank\_cd}: none; fusion changes the hit list for at least 25\% of probes at
  $m = 1{,}000$ & \tbd{} \\
H6 padding (M4) & No latency channel wherever the unpadded AUC was $\ge 0.75$ & \tbd{} \\
\bottomrule
\end{tabular}
\caption{Pre-registered hypotheses of the side-channel experiment, abbreviated; the full rules,
including what falsifies each, are in the released pre-registration.}
\label{tab:side-hyp}
\end{table}

\subsection{Results}\label{sec:side-results}

\tbd{X6 full run (GitHub Actions run 40940385704): guards G1 and G2; per hypothesis the verdict,
the paired and one-shot advantages per store, configuration, $m$ and $p$ with intervals; the
latency cost of M1, M2 and M4 and of separation; any leak (none expected). Figure: paired advantage
against $m$ per configuration.}
